\documentclass[10pt,a4paper]{article}
\usepackage[margin=2cm]{geometry}
\usepackage[T1]{fontenc}
\usepackage{mathptmx}
\usepackage{courier}
\usepackage{microtype}
\usepackage{booktabs}
\usepackage{graphicx}
\usepackage[hidelinks]{hyperref}
\graphicspath{{../results/}}
\setlength{\parindent}{0pt}
\setlength{\parskip}{4pt}
\usepackage[font=small]{caption}

% Numbers and tables below are generated by scripts/pipeline.py (nothing typed by hand).
\input{generated/numbers.tex}

\title{Can Golub weighted voting classify malignant breast samples accurately and robustly?}
\author{}
\date{}

\begin{document}
\maketitle
\vspace{-2.5em}
\emph{Applicability study on the Wisconsin Diagnostic Breast Cancer (WDBC) data. Personal seed: \textbf{\seed}.}
\newline
\emph{Code: https://github.com/LeProDeLaBosse/mlbio-wdbc-golub-applicability ; run \texttt{uv run python scripts/run\_all.py}.}

\section{Data exploration}
The table contains \nsamples{} patients (one row = one patient's fine-needle aspirate) and \nfeatures{} numerical features:
ten nuclear properties (radius, texture, perimeter, area, smoothness, compactness, concavity, concave points, symmetry,
fractal dimension), each summarised as mean, standard error and ``worst'' value. There are \nmal{} malignant (\pctmal\%) and
\nben{} benign samples and no missing values. The model never sees images or patients, only numbers computed after digitisation,
nucleus segmentation and feature extraction on masses that had already been sampled (not a screening population).

To describe the link with the diagnosis I computed, for each raw feature, the AUC of that feature alone and the class means
(Table~\ref{tab:features}, Figure~\ref{fig:auc}). Size and shape-irregularity features are strongly linked to malignancy:
\texttt{\topA} has AUC \topAauc, then \texttt{\topB} (\topBauc) and \texttt{\topC} (\topCauc). Several features are close to
uninformative (\nearhalf; AUC near 0.5). Many features are strongly correlated (Figure~\ref{fig:corr}), notably radius,
perimeter and area. This exploration uses all patients and is descriptive only: it was not used to choose features.

\section{Weighted-voting implementation}
For each training fold I compute, per feature, the class means $\mu_{g,1},\mu_{g,0}$ and standard deviations
$\sigma_{g,1},\sigma_{g,0}$ (malignant = 1, benign = 0; $n-1$ denominator), the weight
$w_g=(\mu_{g,1}-\mu_{g,0})/(\sigma_{g,1}+\sigma_{g,0})$, keep the ten features with largest $|w_g|$ (set $S$), and score a sample
with $V(x)=\sum_{g\in S} w_g\,\bigl(x_g-(\mu_{g,1}+\mu_{g,0})/2\bigr)$. It is predicted malignant if $V\ge 0$. Nothing is
estimated on validation data. Compared with Golub et al.\ (1999), who used 50 genes and a prediction-strength threshold allowing
abstentions, this follows the assignment: ten features, no abstention.

\section{Evaluation (5-fold stratified CV $\times$ 3 repeats, seed \seed)}
\textbf{Example model} (repeat 1, fold 1; Table~\ref{tab:cm}). Weighted voting: AUC \exwvauc, balanced accuracy \exwvba, with
\exwvfn{} of \exwvmal{} malignant samples missed and \exwvfp{} benign sample wrongly called malignant.

\begin{table}[h]\centering\small
\input{generated/tab_confusion.tex}
\caption{Confusion matrices of the example split (the same validation fold for both methods).}\label{tab:cm}
\end{table}

\textbf{Average over the 15 models} (Table~\ref{tab:summary}): weighted voting reaches AUC \wvauc{} and balanced accuracy
\wvba. Specificity is high (\wvspec) but sensitivity is only \wvsens, so about one malignant sample in four is missed. The AUC is much
higher than what the 0 threshold delivers: the score ranks patients well, but the fixed threshold $V=0$ (midpoint of class means)
is poorly placed for this data. I did not test a recalibrated threshold since it is outside the brief; the explanation that skewed,
unequal-variance features shift the midpoint is a hypothesis.

\begin{table}[h]\centering\small
\input{generated/tab_summary.tex}
\caption{Mean $\pm$ SD over the 15 validation folds (same splits for both methods).}\label{tab:summary}
\end{table}

\newpage

\section{Stability of feature selection}
The same \nalways{} features were selected in all 15 training folds (selection frequency 1.00); the other \nnever{} were never
selected (frequency 0.00; full list in Table~\ref{tab:features}). Mean pairwise Jaccard similarity over the 105 pairs is
\textbf{\jaccard}. The selected features are: \selfeatures. Stable prediction does not follow from this: although the feature set
is identical, balanced accuracy still varies between folds (SD \wvbasd), so the variability comes from estimated weights/midpoints
and from which patients fall in each fold, not from which features are chosen. Conversely, a stable identity is partly a consequence
of redundancy: the selected features are mostly size/concavity variables that are correlated with each other and with the unselected
ones, so a different equally good set could exist. Training sets of different folds overlap strongly, so this Jaccard value is
optimistic about stability under truly new data.

\begin{figure}[h]\centering
\includegraphics[width=0.6\linewidth]{fig3_selection_frequency.png}
\caption{Selection frequency of each feature over the 15 training folds.}\label{fig:freq}
\end{figure}

\section{Comparison with logistic regression}
Logistic regression (30 standardised features, scaler fitted on each training fold, liblinear, $C=1$, threshold 0.5, no tuning)
obtains AUC \lrauc{} and balanced accuracy \lrba, sensitivity \lrsens. On the same 15 splits it is better than weighted voting in every
split for both AUC (mean difference WV $-$ LR $=\diffauc$) and balanced accuracy ($\diffba$); specificity is the same, the gap comes
from sensitivity. Caveats: LR uses 30 features against 10 for WV, so the comparison is of two prescribed pipelines, not of two
algorithms at equal complexity; and the 15 folds are not independent (repeated CV on the same patients), so I report the consistent
direction rather than a significance test.

\begin{figure}[h]\centering
\includegraphics[width=0.75\linewidth]{fig4_model_comparison.png}
\caption{Distribution over the 15 folds of AUC and balanced accuracy.}\label{fig:cmp}
\end{figure}

\section{Clinical conclusion}
\textbf{I would not recommend weighted voting, as implemented here, for predicting a patient's diagnosis.} Missing roughly a quarter
of malignant samples is not acceptable for a diagnostic aid, and logistic regression does clearly better with the same data. Even
logistic regression (sensitivity about 0.96) was only validated internally. What the results support: in an FNA-sampled mass, nuclear
size and concavity features carry a strong signal of malignancy, and a simple linear score ranks patients well (AUC 0.96--0.995).
What they do not support: use in screening or asymptomatic people (prevalence here is 37\%, so predictive values would differ),
prognosis or treatment decisions, performance at another centre or with another imaging/segmentation setup, or any causal or
biological claim about the selected features. The data have no site, date or demographic information and we cannot audit how
diagnoses were verified. Further analyses needed: validation on the 54 held-out patients of the original study and on prospective,
multi-centre data; threshold choice driven by a sensitivity target, selected by inner cross-validation; calibration and confidence
intervals (e.g.\ bootstrap); comparison with cytopathologist reading; and evaluation of the effect of segmentation variability.

\section{Reproducibility of the method from the paper}
In this work the computation itself is fully reproducible: a fixed seed, a locked environment (uv) and a SHA256-checked data file give
identical numbers on each run. However, I did not re-run Golub et al.'s own data, so I cannot judge their numerical results; I can
only comment on what had to be decided here. Several details change the outcome and must be stated: the standard deviation
convention, the number of features and how they are chosen, the decision threshold and handling of $V=0$, the absence or use of
prediction strength/abstention, and that all estimation, including feature selection, is done inside each training fold. To make the
original work more reproducible, the authors should have released the code, the exact patient-level data with identifiers for the
splits, the images or at least the segmentation software and parameters, the full list of selected features per training set, the
random seeds and split definitions, and a pre-specified evaluation protocol with confidence intervals, following a reporting
guideline such as TRIPOD, with the independent test set made available.

\section*{References}
\small
\begin{itemize}\setlength{\itemsep}{1pt}
\item Golub TR et al.\ (1999). Molecular classification of cancer: class discovery and class prediction by gene expression monitoring. \emph{Science} 286(5439):531--537. DOI 10.1126/science.286.5439.531.
\item Wolberg W, Mangasarian O, Street N, Street W. Breast Cancer Wisconsin (Diagnostic). UCI Machine Learning Repository. DOI 10.24432/C5DW2B.
\item Street WN, Wolberg WH, Mangasarian OL (1993). Nuclear feature extraction for breast tumor diagnosis. \emph{IS\&T/SPIE Electronic Imaging} 1905:861--870. DOI 10.1117/12.148698.
\item Wolberg WH, Street WN, Mangasarian OL (1994). Machine learning techniques to diagnose breast cancer from image-processed nuclear features of fine needle aspirates. \emph{Cancer Letters} 77(2--3):163--171. DOI 10.1016/0304-3835(94)90099-X.
\item Pedregosa F et al.\ (2011). Scikit-learn: machine learning in Python. \emph{JMLR} 12:2825--2830. (No other code was reused.)
\end{itemize}
\normalsize

\clearpage
\begin{table}[p]\centering\scriptsize
\input{generated/tab_features.tex}
\caption{All features: class means, univariate AUC (all patients, descriptive) and selection frequency by weighted voting.}\label{tab:features}
\end{table}

\begin{figure}[p]\centering
\includegraphics[width=0.7\linewidth]{fig1_univariate_auc.png}
\caption{Univariate AUC of each feature (all patients, descriptive).}\label{fig:auc}
\end{figure}

\begin{figure}[p]\centering
\includegraphics[width=0.75\linewidth]{fig2_correlation.png}
\caption{Correlation matrix of the 30 features.}\label{fig:corr}
\end{figure}

\end{document}
